\uuid{hbWO}
\titre{Raisonnements}
\auteur{bodin ; croizat ; sacre}
\organisation{exo7}
\licence{CC BY-NC-SA 4.0}
\datecreate{2022-03-08}
\niveau{L1}
\module{Algèbre}
\chapitre{Logique, ensemble, raisonnement}
\sousChapitre{Absurde et contraposée}
\difficulte{3}
\theme{Raisonnements}
\exo7qcmid{hilisit.logique.053}

\contenu{
\begin{question}
\qclassification{100.03}
Pour démontrer que, pour tout $x\in\Rr$, on a $|x-2| \le x^2-x+2$.
\begin{answers}
	\bad{Je distingue les cas $x\ge0$ et $x<0$.}
	\good{Je distingue les cas $x\ge2$ et $x<2$.}
	\bad{Je suis amené à vérifier $x^2-x+2 \ge 0$.}
	\good{Je suis amené à vérifier $x^2-2x+4 \ge 0$.}
\end{answers}
\begin{explanations}
Soit $x\in\Rr$. On discute selon le signe de $x-2$.\\
Si $x-2\ge0$ alors $|x-2| = x-2$ donc $(x^2-x+2) - |x-2| = x^2-x+2 - (x-2) = x^2-2x+4$ toujours positif (car le discriminant est négatif).\\
Si $x-2<0$ alors $|x-2| = -(x-2)$ donc $(x^2-x+2) - |x-2| = x^2-x+2 + (x-2) = x^2$ toujours positif.
\end{explanations}
\end{question}
}
